%---------------------------------------------------------------------------------------------------------------------------- 
\graphicspath{{Parti/P03-Meccanica_SR/C11-TLV/Figure/}}
\chapter{La triade meccanica: congruenza, equilibrio e principio dei lavori virtuali  }\label{cap:TLV}
%\chaptermark{Classificazione e soluzione dei sistemi di corpi rigidi }
%---------------------------------------------------------------------------------------------------------------------------- 


\begin{sommario}
	Si formalizza la struttura triadica della meccanica
	delle strutture rigide: congruenza, equilibrio e
	principio dei lavori virtuali sono tre enunciati
	equivalenti, nel senso che due qualsiasi di essi
	implicano il terzo. Assunto il dualismo
	$\bm{B} = \bm{A}^{\top}$ --- stabilito per ispezione
	diretta nel capitolo precedente --- si deriva
	l'\emph{identità bilineare}
	$\bm{r}^{\top}\bm{s} + \bm{u}^{\top}\bm{f} = 0$,
	valida per ogni coppia congruente $(\bm{u}, \bm{s})$
	e ogni coppia equilibrata $(\bm{r}, \bm{f})$.
	Dalla triade si ricavano quattro corollari operativi:
	la condizione di compatibilità statica
	($\bm{U}_q^{\top}\bm{f} = \bm{0}$), il \emph{teorema
		degli spostamenti virtuali} (TSV) --- noto nella
	tradizione come \emph{principio di Müller-Breslau}
	--- per il calcolo delle reazioni, la condizione di
	compatibilità cinematica
	($\bm{R}_\chi^{\top}\bm{s} = \bm{0}$) e il
	\emph{teorema delle forze virtuali} (TFV) --- noto
	nella tradizione come \emph{metodo del carico unitario}
	--- per il calcolo delle componenti di spostamento.
	Si mostra infine come il TFV si estenda al calcolo di
	qualsiasi grandezza cinematica lineare --- spostamento
	assoluto, spostamento relativo, rotazione relativa ---
	mediante la scelta appropriata del problema virtuale
	statico.
\end{sommario}




\section{Il PLV come identità bilineare e il dualismo \texorpdfstring{$\bm{B} = \bm{A}^{\top}$}{B = A\textasciicircum T}}
\label{sec:dualismo}

Il dualismo $\bm{B} = \bm{A}^{\top}$, riconosciuto per
ispezione diretta nel \PARref{sec:prob_stat} del
\CAPref{cap:problemi_cin_stat_SCR}, permette di collegare
il problema cinematico e quello statico attraverso il \PLV. Sia $(\bm{u}, \bm{s})$ una
coppia congruente, con $\bm{A}\,\bm{u} = \bm{s}$, e sia
$(\bm{r}, \bm{f})$ una coppia equilibrata, con
$\bm{A}^{\top}\bm{r} = -\bm{f}$. Il lavoro virtuale
delle reazioni per la configurazione $\bm{u}$ è:
\begin{equation}
	\delta\mathcal{W}_{\mathrm{rea}} =
	\bm{r}^{\top}\bm{s} =
	\bm{r}^{\top}(\bm{A}\,\bm{u}) =
	(\bm{A}^{\top}\bm{r})^{\top}\bm{u} =
	-\bm{f}^{\top}\bm{u} =
	-\delta\mathcal{W}_{\mathrm{att}},
\end{equation}
dove $\delta\mathcal{W}_{\mathrm{att}} = \bm{f}^{\top}\bm{u}$
è il lavoro virtuale delle forze attive. Il lavoro delle
reazioni e quello delle forze attive si bilanciano: si
ottiene così l'\emph{identità bilineare}\index{identità bilineare}:
\begin{equation}\label{eq:id_bil_mec}
	\boxed{\bm{r}^{\top}\bm{s} + \bm{u}^{\top}\bm{f} = 0,}
\end{equation}
valida ogni volta che $(\bm{u}, \bm{s})$ è congruente e
$(\bm{r}, \bm{f})$ è equilibrata. I tre enunciati che
strutturano il problema della meccanica delle strutture
rigide --- congruenza, equilibrio e PLV --- formano la
\emph{triade meccanica}\index{triade meccanica|textbf}:
\begin{enumerate}[label=(\roman*), noitemsep]
	\item \emph{congruenza}: $\bm{A}\,\bm{u} = \bm{s}$;
	\item \emph{equilibrio}: $\bm{A}^{\top}\bm{r} = -\bm{f}$;
	\item \emph{PLV}: $\bm{r}^{\top}\bm{s} +
	\bm{u}^{\top}\bm{f} = 0$.
\end{enumerate}
Due qualsiasi dei tre enunciati implicano il terzo. In
particolare, per cedimenti nulli ($\bm{s} = \bm{0}$),
l'identità~\eqref{eq:id_bil_mec} si riduce a:
\begin{equation}
	\delta\mathcal{W} =
	\bm{f}^{\top}\bm{u} = 0
	\qquad
	\forall\,\bm{u} \in \Ker(\Abb):
\end{equation}
le forze attive non compiono lavoro in alcuna
configurazione cinematicamente ammissibile con vincoli
perfetti. Questo è il contenuto classico del \PLV\ per
i sistemi di corpi rigidi.



%---------------------------------------------------------------------------------------------------------------------------- 
\section{Applicazioni della triade}	\label{sec:applicazioni_triade}
%---------------------------------------------------------------------------------------------------------------------------- 

I tre vertici della triade --- congruenza, equilibrio e PLV
--- implicano ciascuno gli altri due, come si mostrerà nelle
due sottosezioni seguenti. Fissato il PLV nella forma
dell'identità bilineare:
\begin{equation}\label{eq:id_bil_mec_app}
	\bm{r}^{\top}\bm{s} + \bm{u}^{\top}\bm{f} = 0,
\end{equation}
valida ogni volta che $\bm{A}\,\bm{u} = \bm{s}$ e
$\bm{A}^{\top}\bm{r} = -\bm{f}$, si ottengono due famiglie
di risultati a seconda di quale dei due vertici restanti
costituisce il problema reale e quale il problema virtuale.
La struttura della soluzione introdotta nel
\PARref{sec:struttura_soluzione} del \CAPref{cap:problemi_cin_stat_SCR} permette ora di rendere
esplicite le condizioni di compatibilità e di mostrare che esse sono una conseguenza diretta del dualismo
$\bm{B} = \bm{A}^{\top}$.

% ── Congruenza + PLV ─────────────────────────────────────
\subsection{Congruenza e PLV implicano equilibrio}
\label{sec:con_PLV}

Si considera come problema reale il problema statico
$\bm{A}^{\top}\bm{r} = -\bm{f}$ e si applica il
PLV~\eqref{eq:id_bil_mec} con un problema virtuale
cinematico $(\delta\bm{u},\,\delta\bm{s})$ arbitrario
tale che $\bm{A}\,\delta\bm{u} = \delta\bm{s}$.
Sostituendo nella~\eqref{eq:id_bil_mec}:
\begin{equation}\label{eq:id_bil_stat}
	\delta\bm{s}^{\top}\bm{r} +
	\delta\bm{u}^{\top}\bm{f} = 0
	\qquad
	\forall\,(\delta\bm{u},\,\delta\bm{s})
	\text{ congruenti}.
\end{equation}

\begin{corollario}[Compatibilità statica]\index{compatibilita@compatibilità!statica|textbf}
	\label{cor:compat_stat}
	Il problema statico $\bm{A}^{\top}\bm{r} = -\bm{f}$
	ammette soluzione se e solo se:
	\begin{equation}\label{eq:compat_stat}
		\bm{U}_q^{\top}\bm{f} = \bm{0},
	\end{equation}
	ovvero le forze attive $\bm{f}$ compiono lavoro virtuale
	nullo su ogni meccanismo $\delta\bm{u}_q \in \Ker(\Abb)$.
\end{corollario}

\begin{proof}
	Il problema statico ammette soluzione se e solo se
	$\bm{f} \in \Imm(\Abb^{\top})$, ovvero $\bm{f} \perp
	\Ker(\Abb)$ (\APPref{app:applicazionilin_dualita},
	\EQUref{eq:ortogonalita}). Applicando
	la~\eqref{eq:id_bil_stat} con $\delta\bm{s} = \bm{0}$
	e $\delta\bm{u} = \bm{U}_q\,\delta\bm{q}$ per ogni
	$\delta\bm{q} \in \mathbb{R}^{g_\ell}$:
	\begin{equation*}
		\delta\bm{q}^{\top}\bm{U}_q^{\top}\bm{f} = 0
		\qquad \forall\,\delta\bm{q},
	\end{equation*}
	da cui $\bm{U}_q^{\top}\bm{f} = \bm{0}$.
\end{proof}

\begin{oss}
	La condizione di compatibilità statica
	$\bm{U}_q^{\top}\bm{f} = \bm{0}$ mostra che la
	matrice che governa la compatibilità è la trasposta
	della base dei meccanismi: ogni riga di
	$\bm{U}_q^{\top}$ è un meccanismo e impone che le
	forze attive non compiano lavoro in quella direzione.
	Un sistema labile è in equilibrio solo se le forze
	attive sono ortogonali a tutti i meccanismi.
\end{oss}

\begin{corollario}[Teorema degli spostamenti virtuali (TSV)]\index{teorema!degli spostamenti virtuali|textbf}\index{TSV|see{teorema degli spostamenti virtuali}}\index{principio di Müller-Breslau|see{teorema degli spostamenti virtuali}}
	\label{cor:TSV}
	\label{cor:cedimento}
	Si supponga che $\bm{f}$ soddisfi
	la~\eqref{eq:compat_stat}. Sia $\bm{e}_k \in
	\mathbb{R}^m$ il $k$-esimo versore della base canonica.
	Si consideri il \emph{problema virtuale} cinematico
	$\bm{A}\,\bm{u}^*_k = \bm{s}^*$ con
	$\bm{s}^* = \bm{e}_k$. La $k$-esima reazione
	$R_k$ si ottiene dalla soluzione $\bm{u}^*_k$:
	\begin{equation}\label{eq:mueller_breslau}
		R_k = -\bm{f}^{\top}\bm{u}^*_k.
	\end{equation}
\end{corollario}

\begin{proof}
	Si applica la~\eqref{eq:id_bil_stat} alla coppia reale
	$(\bm{r},\,\bm{f})$ e alla coppia virtuale
	$(\bm{u}^*_k,\,\delta\bm{s}^* = \bm{e}_k)$:
	\begin{equation}
		\bm{r}^{\top}\bm{e}_k +
		\bm{f}^{\top}\bm{u}^*_k = 0.
	\end{equation}
	Il primo termine è per definizione $R_k$; riordinando:
	\begin{equation}
		R_k = -\bm{f}^{\top}\bm{u}^*_k. 
	\end{equation}
\end{proof}

\begin{oss}
	Il corollario trasforma il calcolo di una singola
	reazione in un problema cinematico: si impone un
	cedimento unitario al vincolo $k$ e si calcola il
	lavoro delle forze attive nella configurazione che
	ne risulta. Questo è il contenuto del \emph{teorema
		degli spostamenti virtuali} (TSV), noto nella
	tradizione come \emph{principio di Müller-Breslau};
	è alla base della costruzione delle \emph{linee di
		influenza} per le strutture isostatiche, sviluppate
	nel \CAPref{cap:TLV_travi}.
\end{oss}

\begin{oss}[Forma operativa del TSV]
	\label{oss:TSV_operativa}
	La scelta del cedimento virtuale $\bm{s}^* = \bm{e}_k$
	usata nella derivazione è quella più diretta dal punto
	di vista matriciale, ma non quella più comoda nel
	calcolo. Assegnando il cedimento in verso \emph{opposto}
	al verso positivo della reazione --- ovvero $\bm{s}^* =
	-\bm{e}_k$ --- l'identità bilineare~\eqref{eq:id_bil_stat}
	dà direttamente
	\begin{equation}\label{eq:mueller_breslau_op}
		R_k = +\bm{f}^{\top}\bm{u}^*_k,
	\end{equation}
	senza cambio di segno. In questa forma $R_k$ coincide
	con il lavoro virtuale delle forze attive negli
	spostamenti del meccanismo, ed è la scrittura che
	corrisponde al procedimento classico --- \emph{cedere}
	il vincolo $k$-esimo di una quantità unitaria nel verso
	opposto alla reazione assunta positiva, e calcolare il
	lavoro delle forze attive nel conseguente moto rigido.
	Le due forme sono naturalmente equivalenti; la
	specializzazione ai sistemi di travi nel
	\CAPref{cap:TLV_travi} adotta sistematicamente
	la~\eqref{eq:mueller_breslau_op}.
\end{oss}




% ── Equilibrio + PLV ─────────────────────────────────────
\subsection{Equilibrio e PLV implicano congruenza}
\label{sec:eq_PLV}

Si considera come problema reale il problema cinematico
$\bm{A}\,\bm{u} = \bm{s}$ e si applica il
PLV~\eqref{eq:id_bil_mec} con un problema virtuale
statico $(\delta\bm{r},\,\delta\bm{f})$ arbitrario
tale che $\bm{A}^{\top}\delta\bm{r} = -\delta\bm{f}$.
Sostituendo nella~\eqref{eq:id_bil_mec}:
\begin{equation}\label{eq:id_bil_cin}
	\delta\bm{r}^{\top}\bm{s} +
	\delta\bm{f}^{\top}\bm{u} = 0
	\qquad
	\forall\,(\delta\bm{r},\,\delta\bm{f})
	\text{ equilibrati}.
\end{equation}

\begin{corollario}[Compatibilità cinematica]\index{compatibilita@compatibilità!cinematica|textbf}
	\label{cor:compat_cin}
	Il problema cinematico $\bm{A}\,\bm{u} = \bm{s}$
	ammette soluzione se e solo se:
	\begin{equation}\label{eq:compat_cin}
		\bm{R}_\chi^{\top}\bm{s} = \bm{0},
	\end{equation}
	ovvero i cedimenti $\bm{s}$ compiono lavoro virtuale
	nullo su ogni stato di autoequilibrio
	$\delta\bm{r} \in \Ker(\Abb^{\top})$.
\end{corollario}

\begin{proof}
	Il problema cinematico ammette soluzione se e solo se
	$\bm{s} \in \Imm(\Abb)$, ovvero $\bm{s} \perp
	\Ker(\Abb^{\top})$ (Lezione~04). Applicando
	la~\eqref{eq:id_bil_cin} con $\delta\bm{f} = \bm{0}$
	e $\delta\bm{r} = \bm{R}_\chi\,\delta\bm{\chi}$ per ogni
	$\delta\bm{\chi} \in \mathbb{R}^{g_{\mathscr{i}}}$:
	\begin{equation}
		\delta\bm{\chi}^{\top}\bm{R}_\chi^{\top}\bm{s} = 0
		\qquad \forall\,\delta\bm{\chi},
	\end{equation}
	da cui $\bm{R}_\chi^{\top}\bm{s} = \bm{0}$.
\end{proof}

\begin{oss}
	La condizione di compatibilità cinematica
	$\bm{R}_\chi^{\top}\bm{s} = \bm{0}$ mostra che la
	matrice che governa la compatibilità è la trasposta
	della base degli stati di autoequilibrio: ogni riga
	di $\bm{R}_\chi^{\top}$ è uno stato di autoequilibrio
	e impone che i cedimenti non compiano lavoro in quella
	direzione. Un sistema iperstatico ammette una
	configurazione compatibile con i cedimenti solo se
	questi sono ortogonali a tutti gli stati di
	autoequilibrio.
\end{oss}

\begin{corollario}[Teorema delle forze virtuali (TFV)]\index{teorema!delle forze virtuali|textbf}\index{TFV|see{teorema delle forze virtuali}}\index{metodo!del carico unitario|see{teorema delle forze virtuali}}
	\label{cor:carico_unitario}
	Si supponga che $\bm{s}$ soddisfi
	la~\eqref{eq:compat_cin}. Sia $\bm{e}_j \in
	\mathbb{R}^n$ il $j$-esimo versore della base canonica.
	Si consideri il \emph{problema virtuale} statico
	$\bm{A}^{\top}\bm{r}^*_j = -\bm{f}^*$ con
	$\bm{f}^* = \bm{e}_j$. La $j$-esima componente
	$u_j$ si ottiene dalla soluzione $\bm{r}^*_j$:
	\begin{equation}\label{eq:carico_unitario}
		u_j = -\bm{s}^{\top}\bm{r}^*_j.
	\end{equation}
\end{corollario}

\begin{proof}
	Si applica la~\eqref{eq:id_bil_cin} alla coppia reale
	$(\bm{u},\,\bm{s})$ e alla coppia virtuale
	$(\bm{r}^*_j,\,\bm{f}^* = \bm{e}_j)$:
	\begin{equation}
		\bm{s}^{\top}\bm{r}^*_j +
		\bm{u}^{\top}\bm{e}_j = 0.
	\end{equation}
	Il secondo termine è per definizione $u_j$;
	riordinando:
	\begin{equation}
		u_j = -\bm{s}^{\top}\bm{r}^*_j. 
	\end{equation}
\end{proof}

\begin{oss}
	Il corollario è il duale esatto del precedente:
	il calcolo di una componente di spostamento si
	riconduce a un problema statico, applicando una
	forza unitaria nella direzione duale e calcolando
	il lavoro delle reazioni nei cedimenti assegnati.
	Questo è il contenuto del \emph{teorema delle forze
		virtuali} (TFV), noto nella tradizione come
	\emph{metodo del carico unitario}. La specializzazione
	ai sistemi di travi --- in cui l'identità bilineare si
	arricchisce del contributo delle distorsioni nelle
	sezioni interne --- è sviluppata nel
	\CAPref{cap:TLV_travi}, con applicazioni a sistemi
	isostatici soggetti a cedimenti e distorsioni
	assegnati. Nella versione per sistemi deformabili,
	il TFV è inoltre alla base del \emph{metodo delle
		forze} (Vol.~II).
\end{oss}


\begin{oss}[Il dualismo nelle matrici di compatibilità]
	I quattro corollari mostrano che le condizioni di
	compatibilità non sono indipendenti dalla struttura
	della soluzione: la matrice di compatibilità del
	problema statico è $\bm{U}_q^{\top}$, trasposta
	della base dei meccanismi, e quella del problema
	cinematico è $\bm{R}_\chi^{\top}$, trasposta della
	base degli stati di autoequilibrio. Questo non è una
	coincidenza: è la manifestazione diretta del dualismo
	$\bm{B} = \bm{A}^{\top}$. La stessa geometria di
	vincolo che genera i meccanismi (nucleo di $\Abb$)
	determina anche le condizioni che le forze attive
	devono soddisfare per essere equilibrabili, e
	viceversa. Cinematica e statica sono governate dalla
	stessa matrice, letta in direzioni opposte.
\end{oss}


\begin{oss}[Spostamento generico, relativo e rotazione]
	Il Corollario~\ref{cor:carico_unitario} fornisce
	direttamente le componenti $u_j$ del vettore di
	configurazione, associate ai poli $O_i$. La sua
	portata è però più generale, poiché la scelta del
	problema virtuale statico è arbitraria.
	
	\medskip
	\noindent\textit{Spostamento assoluto.}
	Per calcolare lo spostamento $\eta_k$ di un punto
	generico $P_k$ appartenente al corpo $\mathcal{C}_i$
	secondo una direzione $\ab$, si costruisce il
	problema virtuale con una forza fittizia unitaria
	$F_k^* = 1$ applicata in $P_k$ nella direzione
	$\ab$. La soluzione $\bm{r}^*_k$ del problema
	$\bm{B}\,\bm{r}^*_k = -\bm{f}^*$ fornisce:
	\begin{equation}\label{eq:carico_unitario_gen}
		\eta_k = -\bm{s}^{\top}\bm{r}^*_k.
	\end{equation}
	
	\medskip
	\noindent\textit{Spostamento relativo.}
	Per calcolare lo spostamento relativo
	$\Delta\eta = \eta_i - \eta_j$ tra due punti $P_i$
	e $P_j$ appartenenti a corpi distinti secondo la
	direzione della congiungente $P_iP_j$, si applica
	una coppia di forze fittizie unitarie antagoniste:
	$F^* = 1$ in $P_i$ e $F^* = -1$ in $P_j$, entrambe
	nella direzione $\ab = \overline{P_iP_j}/
	\lvert\overline{P_iP_j}\rvert$. Il vettore $\bm{f}^*$
	risultante è la sovrapposizione dei due contributi,
	e la soluzione $\bm{r}^*$ fornisce:
	\begin{equation}\label{eq:scorrimento_relativo}
		\Delta\eta = -\bm{s}^{\top}\bm{r}^*.
	\end{equation}
	
	\medskip
	\noindent\textit{Rotazione relativa.}
	Per calcolare la rotazione relativa
	$\Delta\vartheta = \vartheta_j - \vartheta_i$
	tra due corpi $\mathcal{C}_j$ e $\mathcal{C}_i$,
	si applica una coppia di momenti fittizi unitari
	antagonisti: $\mu^* = +1$ su $\mathcal{C}_j$ e
	$\mu^* = -1$ su $\mathcal{C}_i$. La soluzione
	$\bm{r}^*$ del corrispondente problema virtuale
	fornisce:
	\begin{equation}\label{eq:rotazione_relativa}
		\Delta\vartheta = -\bm{s}^{\top}\bm{r}^*.
	\end{equation}
	
	\medskip\noindent
	In tutti e tre i casi la struttura è identica: la
	quantità cinematica cercata è la duale del sistema
	di forze fittizie applicato nel problema virtuale.
	Il metodo del carico unitario si applica a qualsiasi
	grandezza cinematica che sia un funzionale lineare
	di $\bm{u}$, ovvero della forma $\bm{c}^{\top}\bm{u}$:
	è sufficiente costruire il problema virtuale con il
	sistema di forze che ha $\bm{c}$ come vettore delle
	forze generalizzate.
\end{oss}

%----------------------------------------------------------------------------------------------------------------------------
\section{Esempi applicativi}
\label{sec:esempi_PLV}
%----------------------------------------------------------------------------------------------------------------------------

Gli esempi seguenti riprendono i sistemi già studiati
nel \CAPref{cap:problemi_cin_stat_SCR}, dove sono stati
costruiti la matrice dei vincoli $\bm{A}$, classificato
il sistema e risolti i problemi cinematico e statico.
Qui si applicano i corollari della triade per calcolare
reazioni e spostamenti tramite problemi virtuali.


\begin{esempio}[Trave semplicemente appoggiata ---
	applicazione della triade]
	\label{ex:trave_semplice_PLV}
	Si riprende la trave semplicemente appoggiata
	dell'\ESEref{ex:trave_semplice}: cerniera in $A$,
	carrello in $B = (\ell, 0)$, sistema isostatico con
	cedimento $\bar{v}_B$ al carrello $B$ e carico $F$
	applicato in $x = a$. Calcolare la reazione $R_B$
	tramite il principio di Müller-Breslau e lo spostamento
	$v(a)$ tramite il metodo del carico unitario.


\begin{svolgimento}
	
	\esottotit{Applicazione del principio di Müller-Breslau.}\index{teorema!degli spostamenti virtuali}
	Si assegna cedimento unitario al vincolo~3
	(carrello $B$):
	\begin{equation*}
		\bm{s}^* = \bm{e}_3 =
		\Bigl\{0\; 0\; 1\Bigr\}^{\!\top}.
	\end{equation*}
	Il problema virtuale cinematico
	$\bm{A}\,\bm{u}^* = \bm{s}^*$ dà:
	\begin{equation*}
		u^* = 0, \qquad v^* = 0, \qquad
		\vartheta^* = \frac{1}{\ell}.
	\end{equation*}
	Il lavoro delle forze attive nel campo virtuale è:
	\begin{equation*}
		\bm{f}^{\top}\bm{u}^* =
		\Bigl\{0\; {-F}\; {-Fa}\Bigr\}
		\Bigl\{0\; 0\; 1/\ell\Bigr\}^{\!\top}
		= -\frac{Fa}{\ell}.
	\end{equation*}
	Quindi, per la~\eqref{eq:mueller_breslau}:
	\begin{equation*}
		R_B = -\bm{f}^{\top}\bm{u}^* = \frac{Fa}{\ell},
	\end{equation*}
	in accordo con il risultato diretto del
	\CAPref{cap:problemi_cin_stat_SCR}.
	
	\esottotit{Applicazione del metodo del carico unitario.}\index{teorema!delle forze virtuali}
	Si vuole calcolare lo spostamento verticale $v(a)$
	del punto di applicazione del carico $F$. Si riconosce
	che:
	\begin{equation*}
		v(a) = v + a\vartheta,
	\end{equation*}
	ovvero $v(a)$ è un funzionale lineare di $\bm{u}$ con
	coefficienti $\bm{c} = \Bigl\{0\; 1\; a\Bigr\}^{\!\top}$.
	Il problema virtuale statico con $\bm{f}^* = \bm{c}$
	è $\bm{B}\,\bm{r}^* = -\bm{f}^*$, da cui:
	\begin{equation*}
		R_B^* = \frac{a}{\ell}, \qquad
		Y_A^* = 1 - \frac{a}{\ell}, \qquad
		X_A^* = 0.
	\end{equation*}
	Con il cedimento reale
	$\bm{s} = \Bigl\{0\; 0\; {-\bar{v}_B}\Bigr\}^{\!\top}$:
	\begin{equation*}
		v(a) = -\bm{s}^{\top}\bm{r}^* =
		\frac{a}{\ell}\,\bar{v}_B.
	\end{equation*}
	
	\begin{oss}
		I due corollari mostrano la simmetria della triade
		in azione: per calcolare $R_B$ si risolve un problema
		cinematico virtuale con cedimento unitario al vincolo
		$B$; per calcolare $v(a)$ si risolve un problema
		statico virtuale con forza unitaria nel punto cercato.
		Le due operazioni sono duali l'una dell'altra.
	\end{oss}
	
\end{svolgimento}
\end{esempio}


\begin{esempio}[Trave di Gerber --- principio di Müller-Breslau]
	\label{ex:gerber_PLV}
	Si riprende la trave di Gerber
	dell'\ESEref{ex:gerber}: sistema isostatico con
	$n_c = 2$, carico $F$ nel punto medio di $\mathcal{C}_2$.
	Calcolare la reazione $R_B$ tramite il principio di
	Müller-Breslau.


\begin{svolgimento}
	
	Si assegna cedimento unitario al vincolo~3
	(carrello $B$): $\bm{s}^* = \bm{e}_3$.
	Il problema virtuale $\bm{A}\,\bm{u}^* = \bm{e}_3$
	si risolve sequenzialmente:
	\begin{align*}
		&\text{righe (i)--(ii):} & u_1^* &= v_1^* = 0;\\
		&\text{riga (iii):} & \vartheta_1^* &= \tfrac{1}{2a};\\
		&\text{riga (iv):} & u_2^* &= 0;\\
		&\text{riga (v):} & v_2^* &= 3a\,\vartheta_1^*
		= \tfrac{3}{2};\\
		&\text{riga (vi):} & \vartheta_2^* &=
		-\tfrac{v_2^*}{2a} = -\tfrac{3}{4a}.
	\end{align*}
	Con le forze attive reali
	$\bm{f} = \Bigl\{0\; 0\; 0\; 0\;
	{-F}\; {-Fa}\Bigr\}^{\!\top}$:
	\begin{equation*}
		\bm{f}^{\top}\bm{u}^* =
		(-F)\cdot\tfrac{3}{2} + (-Fa)\cdot
		\Bigl(-\tfrac{3}{4a}\Bigr)
		= -\tfrac{3F}{4},
	\end{equation*}
	quindi:
	\begin{equation*}
		R_B = -\bm{f}^{\top}\bm{u}^* = \tfrac{3F}{4},
	\end{equation*}
	in accordo con il risultato diretto del
	\CAPref{cap:problemi_cin_stat_SCR}.
	
\end{svolgimento}
\end{esempio}


\begin{esempio}[Sistema labile --- compatibilità statica]
	\label{ex:gerber_labile_PLV}
	Si riprende il sistema labile ottenuto dalla trave di
	Gerber rimuovendo il carrello in $B$: $m = 5$,
	$g_\ell = 1$, meccanismo
	$\bm{U}_q = \bigl(0,\;0,\;1,\;0,\;3a,\;
	-\tfrac{3}{2}\bigr)^{\top}$.
	Verificare la compatibilità statica per il carico
	$\bm{f} = \Bigl\{0\; 0\; 0\; 0\; {-F}\;
	{-Fa}\Bigr\}^{\!\top}$.


\begin{svolgimento}
	
	Per il Corollario~\ref{cor:compat_stat}, il problema
	statico ammette soluzione se e solo se
	$\bm{U}_q^{\top}\bm{f} = \bm{0}$, ovvero:
	\begin{equation*}
		\bm{U}_q^{\top}\bm{f} =
		1\cdot \mu(O_1)
		+ 3a\cdot Y_2
		- \tfrac{3}{2}\cdot \mu(O_2) = 0.
	\end{equation*}
	Per il carico assegnato:
	\begin{equation*}
		\bm{U}_q^{\top}\bm{f} = -\tfrac{3Fa}{2} \neq 0.
	\end{equation*}
	La condizione non è soddisfatta: il sistema labile
	non è in equilibrio sotto questo carico. Fisicamente,
	manca il carrello $B$ che fornirebbe il momento
	necessario a bilanciare il carico su $\mathcal{C}_2$.
	
	\begin{oss}
		La condizione $\bm{U}_q^{\top}\bm{f} = \bm{0}$
		è la riscrittura matriciale dell'equilibrio alla
		rotazione del meccanismo: la stessa equazione che
		uno studente di statica elementare scriverebbe
		prendendo i momenti rispetto ad $A$.
	\end{oss}
	
\end{svolgimento}
\end{esempio}
